\subsection{全纯函数}

本节假定 K = C。

3.3.1. 假设 F 拟完备。设 U 是 E 的开子集，f 是从 U 到 F 的映射。下列条件等价：

\begin{enumerate}
\def\labelenumi{(\roman{enumi})}
\item
  \(f\) 是全纯的；
\item
  \(f\) 是可微的；
\item
  \(f\) 局部有界，并且对任意 \(a \in \mathbf{U}\) 和 \(h \in \mathbf{E}\)，在 \(t \mapsto f(a + th)\) 中 0 的某开邻域上定义的函数 \(\mathbf{C}\) 是全纯的；
\item
  \(f\) 局部有界，并且对 F 上每个连续线性形式 \(u\)，取值于 C 的函数 \(u \circ f\) 是全纯的；
\item
  \(f\) 连续且局部有界，并且 F 的对偶中存在一个完备集 H，使得对每个 \(u \circ f\)，\(u \in \mathbf{H}\) 都是全纯的。
\end{enumerate}

当 E 有限维时（分别地，当 F 是巴拿赫空间时），去掉条件 (iii)、(iv) 或 (v)（分别地，条件 (iv) 或 (v)）中的假设“f 局部有界”，所得的条件 (iii′)、(iv′) 或 (v′)（分别地，(iv′) 或 (v′)）仍与这些条件等价。

3.3.2. 假设 F 拟完备。设 U 是 E 的开子集，\((f_{n})\) 是从 U 到 F 的全纯映射序列，具有下列性质：

\begin{enumerate}
\def\labelenumi{(\Alph{enumi})}
\setcounter{enumi}{22}
\tightlist
\item
  U 的每一点都有一个邻域，使得序列 \((f_{n})\) 在该邻域上一致收敛。
\end{enumerate}

则序列 (f＿n) 的极限 \(f\)\((f_{n})\) 是全纯的，导数序列 \((\mathrm{D}f_{n})\)（取值于拟完备空间 \(\mathcal{L}(\mathrm{E};\mathrm{F})\)）具有性质 (W)，且 Df 是 \((\mathrm{D}f_{n})\) 的极限。

3.3.3. 设 U 是 E 的开子集，f 是从 U 到 F 的全纯映射，并假设 F 拟完备。令 \(\mathbf{R} = (\mathbf{R}_{i}) \in (\mathbf{R}_{+}^{*})^{n}\)，并假设多球 \(\mathbf{B}(\mathbf{R})\) 包含于 U 且 f 在 \(\mathbf{B}(\mathbf{R})\) 上有界。于是，对每个 \(\alpha \in N^{n}\) 以及每个 \(x = (x_{i}) \in \mathbf{B}(\mathbf{R})\)，有：

\[
\Delta^ {\alpha} f (0) (x) = \int_ {0} ^ {1} \dots \int_ {0} ^ {1} f (\mathbf {e} (\theta_ {1}) x _ {1}, \dots , \mathbf {e} (\theta_ {n}) x _ {n}) \mathbf {e} (- \alpha_ {1} \theta_ {1} - \dots - \alpha_ {n} \theta_ {n}) d \theta_ {1} \dots d \theta_ {n}
\]

\((\text{其中 } \mathbf{e}(\theta) = \exp 2\pi i \theta)\)。

此外，设 \(\gamma\) 是 \(F\) 上的连续半范数，\(M\) 是在 \(\|f(x)\|_{\gamma}\) 时 \(\|x_{i}\|=R_{i}\) 的上确界。则对所有 \(\|\Delta^{\alpha}f(0)(x)\|_{\gamma}\leqslant M\) 都有 \(x\in B(R)\)，且 \(\|\Delta^{\alpha}f(0)\|_{\gamma}\leqslant MR^{-\alpha}\)。最后，\(f\) 在 0 处的展开级数的收敛域包含多圆盘 \(B(R)\) 的内部。

3.3.4. 保持 3.3.3 的假设，并进一步假设 \(\mathbf{E}_i = \mathbf{C}\)。设 \(\sum X^{\alpha}c_{\alpha}\) 是 \(f\) 在 0 处的展开级数。有：

\[
c _ {\alpha} = \mathbf {R} ^ {- \alpha} \int_ {0} ^ {1} \dots \int_ {0} ^ {1} f (\mathbf {e} (\theta_ {1}) \mathrm{R} _ {1}, \dots , \mathbf {e} (\theta_ {n}) \mathrm{R} _ {n}) \mathbf {e} (- \alpha_ {1} \theta_ {1} - \dots - \alpha_ {n} \theta_ {n}) d \theta_ {1} \dots d \theta_ {n}
\]

以及：

\[
\| c _ {\alpha} \| _ {\gamma} \leqslant \mathbf {R} ^ {- \alpha} \sup _ {x \in \mathbf {B} (\mathbf {R})} \| f (x) \| _ {\gamma}
\]

（“柯西不等式”）。级数 \(\sum_{\alpha} X^{\alpha} c_{\alpha}\) 的严格收敛域包含 B(R) 的内部。

3.3.5. 假设 E 有限维且 F 拟完备。于是 \(\mathcal{H}(\mathbf{E};\mathbf{F})\) 中存在唯一的级数 \(f_{0}\)，它（对 E 上每个范数）具有无穷收敛半径，并满足对所有 \(f(x)=f_{0}(x)\) 都有 \(x\in E\)。

3.3.6. 若 \(f\) 是从 E 到 F 的全纯映射且 \(f(E)\) 有界，则函数 \(f\) 为常值函数（“刘维尔定理”）。

3.3.7. 假设 \(E \neq 0\)。设 f 是从 E 的开集 U 到 F 的全纯映射。设 \(a\) 是 U 中一点，\(\gamma\) 是 F 上的连续半范数。对包含于 U 的 a 的每个邻域 V，都存在 \(x \in V\)、\(x \neq a\)，使得：

\[
\| f (a) \| _ {\gamma} \leqslant \| f (x) \| _ {\gamma}.
\]

若此外 F = C，且 f 在 a 的任何邻域上都不为常值，则有 \(|f(a)| < \sup_{x \in V, x \neq a} |f(x)|\)，并且映射 f 在 a 的某邻域上是开映射。最后，设 B 是闭包包含于 U 的有界开集，B' 是其

边界。有：

\[
\sup _ {x \in \overline {{\mathbf {B}}}} | f (x) | = \sup _ {x \in \mathbf {B} ^ {\prime}} | f (x) |
\]

（“最大值原理”）。

3.3.8. 假设 E 有限维。设 U 是 E 的开子集，S 是 E 的余维 \(\geqslant2\) 的向量子空间，f 是从 \(\mathrm{U}\cap(\mathrm{E}-\mathrm{S})\) 到 F 的全纯映射。则 f 可连续延拓为从 U 到 F 的全纯映射。

假设 E = C，且 \(0 \leqslant \lambda < \mu \leqslant +\infty\)。设 f 是从开集 \(U = \{z \in C | \lambda < |z| < \mu\}\) 到 F 的全纯映射。存在唯一的 F 中元素族 \((a_{n}(f))_{n \in \mathbb{Z}}\)，使得对 U 的每个紧子集 H，族 \((a_{n}(f)z^{n})_{n \in \mathbb{Z}}\) 一致可求和，当 z 遍历 H 时其和为 \(f(z)\)（“洛朗展开”）。

再假设 \(\lambda=0\)。称 \(f\) 在点 \(x=0\) 处的阶为满足 a＿\(n\)\(a_{n}(f)\neq0\)(\(f\))0 的整数 n 的下确界（有限或无穷）。若存在 0 的邻域 V，使得 f 在 V-\(\{0\}\) 中有界，则 \(f\) 在点 \(x=0\) 处的阶为 0，并可连续延拓为开集 \(|z|<\mu\) 上的全纯函数。设 \(p\) 是整数且 p\(>0\)；若 \(f\)\(-p\) 在点 \(x=0\) 处的阶为 -\(p\)，则称 0 是 \(f\) 的 p 阶极点。

3.3.10. 假设 E = C 且 F 是赋范空间。设 f 是从 E 的开单位圆盘 U 到 F 的全纯映射，满足 \(f(0) = 0\)，并令 \(M = \sup_{z \in U} \|f(z)\|\)。于是对所有 \(\|f(z)\| \leqslant M \cdot |z|\) 都有 \(z \in U\)（“施瓦茨引理”）。

